\chapter{Wykład 14}

Nasza definicja relacji $E$ działa dobrze jeśli nasz pierścień jest dziedziną całkowitości. Jednak gdy pierścień nie jest dziedziną, definicja ta się psuje. Zobaczmy dokładniej -- spróbujemy jej użyć na pierścieniu, który ma dzielniki zera (nie jest dziedziną), sprawdzając przechodniość. Załóżmy, że mamy trzy ułamki:
\begin{enumerate}[label=\arabic*)]
	\item $\frac{r}{s} E \frac{r'}{s'}$ (czyli $rs' - r's = 0$)
	\item $\frac{r'}{s'}E \frac{r''}{s''}$ (czyli $r's'' - r''s' = 0$ )
\end{enumerate}
Chcemy pokazać, że $\frac{r}{s}E \frac{r''}{s''}$, czyli że $rs'' - r''s = 0$. Mnożymy pierwsze równanie przez $s''$ a drugie przez $s$:
 \begin{enumerate}[label=\arabic*)]
	\item $rs's'' - r'ss'' = 0$
	\item $r's''s - r''s's = 0$
\end{enumerate}
Dodając stronami, (zauważmy, że wyraz $r's''s$ się redukuje) otrzymujemy $rs's'' - r''s's = 0$.

Wyciągając  $s'$ przed nawias otrzymujemy $s'\left( rs'' - r''s  \right) = 0$.

I tu napotykamy na ścianę:
\begin{itemize}
	\item W dziedzinie całkowitości: skoro $s' \in S$ (więc $s' \neq 0$), to możemy skrócić równanie przez $s'$. Wtedy  $rs'' - r''s = 0$. Koniec dowodu.
	 \item W naszym przypadku (z dzielnikami zera): $s'$ może być dzielnikiem zera. Fakt, że $s'\cdot \left( \text{coś} \right)  = 0$, nie oznacz, że $\left( \text{coś} \right) = 0$. Zatem bez nowej definicji, relacja nie byłaby przechodnia, czyli nie byłaby relacją równoważności. Konstrukcja upada.
\end{itemize}

Nowa definicja mówi: nie wymagamy aby różnica $rs' - r's$ była zerem, wymagamy aby była ,,unicestwiana'' przez jakiś element z  $S$:
\[
	\left( r,s \right) \sim \left( r', s' \right) \iff \left(\exists t \in S \right)\left( t \left( rs' - r's \right) = 0 \right)   
.\] 

Sprawdźmy teraz przechodniość -- mieliśmy moment krytyczny: $s'\left( rs'' - r''s \right) = 0$. W starej definicji to był koniec. Natomiast w nowej, to równanie mówi nam, że $s'$ jest naszym ,,świadkiem''. Skoro $s' \in S$ i $s'$ pomnożone przez nawias daje $0$, to zgodnie z nową definicją para $\left( r,s \right) $ jest w relacji z $\left( r'',s'' \right) $.

Definicja jest ,,szyta na miarę'', aby uratować przechodniość. Świadek $t$ (lub  $s''$ ) to element, który ,,zabija'' szum informacyjny wynikający z dzielników zera.

 \noindent Intuicja geometryczna (kiełki).\\

 Zobaczmy to na przykładzie kiełków(SPOJLER). Powiedzieliśmy, że $\frac{f}{g} = \frac{h}{k}$ w lokalizacji, jeśli są równe na pewnym otoczeniu $\mathcal{U}$.

 Zapiszmy to algebraicznie:
 \[
 f\left( x \right) k\left( x \right) - h\left( x \right) g\left( x \right) = 0 \text{, dla } x \in \mathcal{U}
 .\] 

 Poza zbiorem $\mathcal{U}$ ta różnica może być ogromna. Równanie globalne nie jest spełnione ($fk - hg \neq 0$). Ale mamy świadka. Niech $t \in S$ będzie funkcją, która
 \begin{itemize}
 	\item Jest równa $1$ na jakimś podzbiorze wewnątrz $\mathcal{U}$.
	\item Jest równa $0$ wszędzie tam, gdzie $fk - hg \neq 0$ (poza $\mathcal{U}$ )
 \end{itemize}

Wtedy globalnie iloczyn $t\left( fk - hg \right) $ będzie tożsamościowo równy zerze na całym $\mathbb{R}$. Funkcja $t$ ,,wycina'' (zeruje) te fragmenty dziedziny, które nas nie interesują (są dalekie od zera), i zostawia wszystko to, co jest blisko.



Można też lokalizować pierścienie które nie są dziedzinami. Niech $S\subseteq R$ (gdzie $S$ jest podzbiorem multiplikatywnym) będzie pierścieniem przemiennym (niekoniecznie dziedzina).

Wtedy relacja $\left( r,s \right) E \left( r',s' \right) \iff rs' = r's$ nie musi być relacją równoważności (przechodniość się psuje).

Definiujemy relację $\sim $ na $R\times S$ poniższym wzorem:
	\[
		\left( r,s \right) \sim \left( r',s' \right) \iff \left(\exists s''\in S \right)\left( s''(rs'-r's) = \mathbf{0} \right)   
	.\] 
Tak zdefiniowana relacja $\sim$ jest relacją równoważności i $R_{S} := R\times \mfaktor{S}{\sim} $ jest pierścieniem.

Zauważmy, że odwzorowanie 

\begin{align*}
	\varphi : R &\longrightarrow R_{S} \\
	r &\longmapsto \varphi (r) = \frac{r}{1} = \mfaktor{\left( r,s \right) }{\sim  } 
,\end{align*}
jest homomorfizmem ale niekończenie $,,1-1''$. Twierdzenie (Jeśli $P$ jest ideałem pierwszym w $R$ ($R$ --dziedzina) to $R_{P} = R_{R\setminus P}$ jest lokalny) wciąż zachodzi, pomimo iż $R$ niekoniecznie jest dziedziną.
 
\noindent \textbf{Przykład}

$C\left( \mathbb{R}  \right) = \{f: \mathbb{R} \to \mathbb{R} : f \text{ jest ciągła}\}$.

Jest to pierścień przemienny z $\mathbf{1}$ ale nie dziedziną, bo dla przykładowych funkcji (patrz poniższe grafika) $f,g \neq 0$ ale $fg = 0$.
\begin{figure}[H]
	\centering
	\includegraphics[width=0.8\textwidth]{"funkcje.png"}
	\caption{}
	\label{fig:}
\end{figure}

Niech $\mathcal{I}_{\mathbf{0}} := \{f \in C\left( \mathbb{R}  \right): f\left( \mathbf{0} \right) = 0\} $. Wtedy 
\[
\mfaktor{C\left( \mathbb{R} \right) }{\mathcal{I}_{0}} \cong \mathbb{R}
.\] 
Dlaczego?



Ponieważ $\mathbb{R}$ jest ciałem więc $\mathcal{I}$ jest ideałem maksymalnym, więc w szczególności jest ideałem pierwszym.
Zatem $C\left( \mathbb{R}  \right)_{\mathbb{R} , \mathbf{0}} := C\left( \mathbb{R}  \right)_{\mathcal{I}_{0}}$ jest lokalizacją.

\noindent $C\left( \mathbb{R}  \right) _{\mathbb{R} , 0}$ jest to pierścień \textbf{kiełków} ciągłych funkcji rzeczywistych w 0.

\noindent Rozważmy \begin{align*}
	\varphi  : C\left( \mathbb{R}  \right)  &\longrightarrow C\left( \mathbb{R}  \right) _{\mathbb{R} , 0} \\
	f &\longmapsto \varphi  (f) = \frac{f}{1}
.\end{align*}

\noindent Tutaj (bardzo nietypowe) $\varphi  $ jest epimorfizmem bo
\[
\left(\forall \tfrac{f}{g}\in C\left( \mathbb{R} \right)_{\mathbb{R},0} \right)\left(\exists \mathcal{U}\ni 0 \right)\left(\exists h \right)\left( h\restriction_{\scriptscriptstyle \mathcal{U}} = \tfrac{f}{g}\restriction_{\scriptscriptstyle \mathcal{U}}  \right) 
.\] 
\noindent Biorąc $t\in C\left( \mathbb{R}  \right) \setminus \mathcal{I}_{0}$, takie że $t\restriction_{\scriptscriptstyle \mathbb{R} \setminus \mathcal{U}}$ (czyli $t$ obcięte do $\mathbb{R}\setminus \mathcal{U}$ ). Mamy $t\left( gh-f \right) = 0$, a więc $\frac{f}{g}= \frac{h}{1}= \varphi \left( h \right)  $

\begin{figure}[H]
	\centering
	\includegraphics[width=0.8\textwidth]{"epikiełki.png"}
	\caption{}
	\label{fig:}
\end{figure}

\noindent \textbf{Komentarz:}\\
\begin{itemize}
	\item Panel 1 (Element lokalny): Wykres przedstawia ułamek $\frac{f}{g} \in C\left( \mathbb{R} \right)_{\mathbb{R}, 0}$, który jest dobrze określony w otoczeniu zera, ale posiada osobliwość w $x = 1.2$. Nie jest to więc funkcja z  $C\left( \mathbb{R} \right) $.
	\item Panel 2 (Narzędzie -- funkcja t): Funkcja typu ,,bump''. Przyjmuje wartość $1$ w centrum, by zachować własność kiełka i spada do $0$ przed osiągnięciem punktu osobliwego.
	\item Panel 3 (Wynik -- funkcja globalna $h$): Iloczyn  $h = t \cdot \frac{f}{g}$. W otoczeniu zera -- $h = \frac{f}{g}$, czyli lokalnie jest zgodna. Natomiast globalnie funkcja $h$ jest ciągła na całym  $\mathbb{R}$. Osobliwość została ,,wyzerowana''.
\end{itemize}
\noindent \textbf{Wniosek:} Każdy element lokalizacji (ułamek) pochodzi od pewnej funkcji globalnej $h$. Oznacza to, że homomorfizm $\varphi: C\left( \mathbb{R} \right) \to C\left( \mathbb{R} \right)_{\mathbb{R}, 0} $ jest epimorfizmem (surjekcją). 


\noindent Podobnie 
 \[
 \left(\forall h_1,h_2 \in C\left( \mathbb{R} \right)\right)\left( \varphi\left( h_1 \right) = \varphi\left( h_2 \right) \iff \left(\exists \mathcal{U} \ni 0 \right)\left( h_1\restriction_{\scriptscriptstyle \mathcal{U}} = h_2\restriction_{\scriptscriptstyle \mathcal{U}}  \right)     \right) 
 .\] 
 \noindent \textbf{Wniosek:} $\varphi  $ \textbf{nie} jest $,,1-1''$.

 Elementy w  $C_{\mathbb{R} ,0}$ są równe $\iff $ są ,,lokalnie równe wokół zera''

\begin{figure}[H]
	\centering
	\includegraphics[width=0.8\textwidth]{"kiełki.png"}
	\caption{}
	\label{fig:}
\end{figure}
Mamy powyżej 3 funkcje:
\begin{itemize}
	\item $f\left( x \right) $- czarna, oryginał (cosinus)
	\item $g\left( x \right) $ - czerwona, zgadza się z $f$ na otoczeniu $U = \left( -1, 1 \right) $ 
	\item $h\left( x \right) $ - niebieska, zgadza się z $f$ na otoczeniu (węższym) $V = \left( -0.5, 0.5 \right) $
\end{itemize}

To ilustruje kluczową własność relacji równoważności kiełków:

Dwie funkcje definiują ten sam kiełek, jeśli istnieje \textbf{jakiekolwiek} otoczenie, na którym są równe. Nie musi to być to samo otoczenie dla każdej pary.

\begin{itemize}
	\item Para $\left( f,g \right) $ zgadza się na dużym otoczeniu $U$.
	\item Para $\left( f,h \right) $ zgadza się na mniejszym otoczeniu $V$.
	\item Ponieważ $h$ zgadza się z $f$ na małym otoczeniu $V$, a $g$ zgadza się na większym otoczeniu $U$, które zawiera  $V$, to wszystkie trzy są tym samym kiełkiem.
\end{itemize}

\textbf{Stąd nazwy:}  
 \begin{enumerate}[label=\arabic*)]
 	\item \textbf{kiełki} , bo istotnie jest zachowanie wokół zera, czyli to co kiełkuje,
	\item \textbf{lokalizacja}, bo istotne jest tylko lokalne zachowanie funkcji.
 \end{enumerate}

 \section{Rozkład w pierścieniu wielomianów.}

 Niech $R$ będzie UFD oraz $K=R_{0}$. \textbf{Cel:} $R\left[ X \right] $ jest UFD.

\noindent Chcemy udowodnić, że jeśli pierścień współczynników $R$ jest dziedziną z jednoznacznością rozkładu (UFD), to pierścień wielomianów nad nim  $R\left[ X \right] $ również nią jest.

\noindent Dlaczego to nie jest oczywiste?\\
\begin{enumerate}[label=\arabic*.]
	\item Pułapka podpierścienia: Wiemy, że $K\left[ X \right] $ (gdzie $K$ to ciało ułamków $R$ ) jest UFD (bo jest dziedziną euklidesową). Ale $R\left[ X \right] $ jest podpierścieniem $K\left[ X \right] $ -- zauważmy że bycie podpierścieniem UFD nie gwarantuje bycia UFD.
	\noindent \textbf{Przykład:} $\mathbb{Z}\left[ \sqrt{-5}  \right] \underset{\scriptscriptstyle\text{podp}}{\subseteq} \mathbb{C}$. Ciało liczb zespolonych jest UFD, ale $\mathbb{Z}\left[ \sqrt{-5}  \right] $ nie jest -- $6 = 2\cdot 3 = \left( 1+\sqrt{-5}  \right)\left( 1-\sqrt{-5}  \right) $. Zatem sam fakt bycia ,,wewnątrz'' nam nie wystarcza.
	\item Różnica w odwracalności: Kluczowa obserwacja: rozważmy wielomian $2X$.
		 \begin{itemize}
			\item W $\mathbb{Q}\left[ X \right] $ (nad ciałem): $2$ jest odwracalna (jednostka). Możemy napisać $2 \cdot X = 2X$. Ponieważ $2$ jest jednostką, to $2X$ jest \textbf{stowarzyszone} z $X$ ($2X \sim X$). Mnożenie przez jedność nie wpływa na rozkładalność, stopień wynosi $1$, więc jest to wielomian nierozkładalny. Wielomian jest nierozkładalny w tym sensie, że nie rozpada się na dwa wielomiany niższego stopnia.
			\item W $\mathbb{Z}\left[ X \right]$ (nad pierścieniem): Liczba $2$ nie jest odwracalna. Zatem zapis  $2X = 2 \cdot X$ jest rozkładem na dwa elementy nieodwracalne (stała $2$ i wielomian $X$ )-- definicja ,,czynnika nierozkładalnego'' w $R\left[ X \right] $ obejmuje też stałe z $R$.
		\end{itemize}
\end{enumerate}
\noindent \textbf{Wniosek:} Arytmetyka w $R\left[ X \right] $ jest ,,bogatsza'' niż w $K\left[ X \right] $, bo dochodzi nam faktoryzacja współczynników. 


 \noindent \textbf{Obserwacja -- różnica w zbiorze jednostek:} Element $c \in R$ może nie być odwracalny w $R$ (nie jest jednostką), stając się jednocześnie elementem odwracalnym (jednością) w  $K$. W konsekwencji,  $c$ może być istotnym czynnikiem w rozkładzie wielomianu w $R\left[ X \right] $, podczas gdy w $K\left[ X \right] $ traktowany jest trywialnie (jako stowarzyszenie).
 \[
  R^{\ast} \subsetneq K^{\ast} \implies \text{ pojęcie nierozkładalności w } R\left[ X \right] \text{ jest bardziej restrykcyjne}
 .\] 

 \noindent \textbf{Idea -- normalizacja do wielomianów pierwotnych:} Dowolny wielomian $f\in K\left[ X \right] $ można przedstawić w postaci iloczynu z $K$ (reprezentującej ,,wspólną część'' współczynników) oraz wielomianu z  $R\left[ X \right] $, którego współczynniki są względnie pierwsze. Strategia polega na sprowadzeniu badania rozkładalności $f$ do badania jego \textbf{części pierwotnej} (wielomian pierwotny -- primitive part). 


\noindent \textbf{Definicja (zbiór reprezentantów):} Niech $P \subseteq R$ będzie ustalonym zbiorem reprezentantów wszystkich klas elementów nierozkładalnych w $R$ względem relacji stowarzyszenia $\sim  $. Oznacza to, że każdy element nierozkładalny $r \in R$ jest stowarzyszony z dokładnie jednym $p \in P$:
\[
\left(\forall r \in R_{\scriptscriptstyle \text{nierozkłdalne}} \right) \left(\exists ! p \in P \right)\left( r \sim p  \right)  
.\] 

\begin{eg}
	Niech $R = \mathbb{Z}$. Naturalnym wyborem na $P$ jest zbiór liczb pierwszych: $\{2,3,5,7,\ldots\} $. Wtedy
	\[
	\left(\forall m \in \mathbb{Z}\setminus \{0, \pm1\}  \right)\left(\exists ! p_1,p_2,\ldots,p_{n} \in P \right)\text{ takie że} \left(\exists ! v_1,\ldots,v_{n} \in \N_{>0} \right)\left( m = \pm p_1^{v_1}\cdot \ldots \cdot p_{n}^{v_{n} } \right)  
	.\] 
Można teraz zapisać (dopuszczając nieskończenie wiele $1$, czyli $p_{i}^{0}$ )
\[
\mathbb{Z} \setminus \{0,\pm 1\} \ni m = \pm \prod_{p \in P} p ^{v_{p}} \text{ dla jedynych $v_{p}\in \N$}
.\] 
\end{eg} 
\begin{definition}
	Funkcję $v_{p} : R\setminus \{0\} \to \mathbb{Z}$ przyporządkowującą elementowi $a$, wykładnik w którym  $p$ występuje w rozkładzie $a$, nazywamy \textbf{waluacją p--adyczna}.
\end{definition} 
\noindent Ogólniej:
\begin{theorem}[Postać normalna UFD]
	Niech $P$ będzie ustalonym zbiorem reprezentantów. Każdy element $a \in R\setminus \{0\} $ posiada \textbf{jednoznaczne} przedstawienie w postaci
	\[
	a = u\cdot \prod_{p \in P} p^{v_{p}\left( a \right) }  
	,\]
	gdzie:
	\begin{enumerate}[label=\arabic*)]
		\item $u \in R^{\ast} $ jest \textbf{jednoznacznie}  wyznaczonym elementem odwracalnym $\left(\exists ! u \in R^{\ast}  \right) $,
		\item $v_{p}\left( a \right) \in \mathbb{Z}$ to wykładnik w którym $p$ występuje w rozkładzie $a$,
		 \item $v_{p}\left( a \right) = 0 $ dla prawie wszystkich (tzn. zbiór $\{p \in P: v_{p}\left( a \right) > 0\} $ jest skończony)
	\end{enumerate}
\end{theorem} 

\noindent Zauważmy, że dla $a \in R\setminus \{0\} $, $v_{p}$ przyjmuje wartości naturalne, a dla $a \in K\setminus \{0\} $ (czyli ciała), $v_{p}$ przyjmuje wartości całkowite (wykładnik może być ujemny).

\begin{eg}
	Rozważmy ułamek $a = \frac{-7}{24}$. Wtedy $a = \left( -1 \right) 2^{-3}3^{-1}7^{1}$ oraz
	\begin{itemize}
		\item $v_{2}\left( a \right) = -3$,
		\item $v_{3}\left( a \right) = -1$,
		\item $v_{7}\left( a \right) = 1$,
		\item $v_{p}\left( a \right) = 0$, dla $p \neq 2,3,7$
	\end{itemize}
\end{eg} 
\begin{lemma}[Wprost z definicji UFD.]
	Niech $a,b \in K\setminus \{\mathbf{0}\} $.
	\begin{enumerate}[label=\arabic*)]
		\item $a\in R \iff \left(\forall p \in P \right)\left( v_{p}(a) \ge 0 \right)  $ 
		\item $\left(\forall p \in P \right)\left( v_{p}\left( ab \right) = v_{p}(a) + v_{p}(b) \right)  $
	\end{enumerate}
\end{lemma} 
\begin{proof}
	TODO
\end{proof} 
\begin{definition}
	Dla $f = a_0 + a_1X + a_2X^2 + \ldots + a_{n} X^{n} \in K\left[ X \right] \setminus \{\mathbf{0}\} $ oraz $p \in P$:
	\begin{enumerate}[label=\arabic*)]
		\item $v_{p}\left( f \right) := \min \{v_{p}\left( a_0 \right) , \ldots, v_{p}\left( a_{n}  \right) \} $,
		\item $\cont\left( f \right) := \prod_{p\in P} p^{v_{p}(f)} \in K\setminus \{0\}  $
	\end{enumerate}
	$\cont$ to tłumaczymy na ,,zawartość''. Jest to funkcja:
	\[
	\cont: K\left[ X \right] \setminus \{\mathbf{0}\} \to K\setminus \{\mathbf{0}\} 
	.\] 
\end{definition} 
\begin{remark}
	Ważne:
	\begin{enumerate}[label=\roman*)]
		\item $f\in \mathbb{Z} \left[ X \right] \setminus \{\mathbf{0}\} \implies \cont\left( f \right) = \text{NWD}\left( a_0,\ldots,a_{n}  \right) .$
		\item $\cont\left( f \right) \in R^{\ast} \iff \cont\left( f \right) = 1$
		\item $\left(\forall a \in K\left[ X \right] \setminus \{\mathbf{0}\}  \right)\left(\exists \varepsilon \in R^{\ast}  \right)\left( a = \varepsilon \cont \left( a \right)  \right)   $.
		\item $\left(\forall a \in K\left[ X \right] \setminus \{\mathbf{0}\}  \right)\left( \cont \left( \cont \left( a \right)  \right) = \cont \left( a \right)  \right)  $
		\item $\left(\forall a,b \in K\setminus \{\mathbf{0}\}  \right)\left( \cont \left( ab \right) = \cont \left( a \right) \cont \left( b \right)  \right)  $.
	\end{enumerate}
\end{remark} 
\begin{eg}
	Niech $R = \mathbb{Z} $.
	\begin{itemize}
		\item $v_2\left( 8X + \frac{2}{3} \right) = \min \{v_2\left( 8 \right) , v_2\left( \frac{2}{3} \right) \} = \min \{3,1\} = 1$
		\item $v_2\left( 8X + \frac{2}{3} \right) = \min \{v_3\left( 8 \right) , v_3\left( \frac{2}{3} \right) \} = \min \{0, -1\} = -1$
	\end{itemize}
	Stąd $\cont \left( 8X + \frac{2}{3} \right) = 2^{1}\cdot 3^{-1} = \frac{2}{3}$ 
	
	Zauważmy: 
	\[
		\frac{8X + \frac{2}{3}}{\cont \left( 8X + \frac{2}{3} \right) } = 12X + 1 
	,\]
	a $\cont \left( 12X + 1 \right) = 1$
\end{eg} 
\begin{lemma}
	Niech $f\in K\left[ X \right] \setminus \{\mathbf{0}\}, a \in K\setminus \{\mathbf{0}\}  $
	\begin{enumerate}[label=\arabic*)]
		\item $f \in R\left[ X \right] \iff \cont \left( f \right) \in R$,
		\item $\cont \left( af \right) = \cont \left( a \right) \cont \left( f \right) $,
		\item $\cont \left( \frac{f}{\cont \left( f \right) } \right) = 1$.
	\end{enumerate}
\end{lemma} 
\begin{proof}
	TODO
\end{proof} 
\begin{lemma}[Guassa]
	\[
	\left(\forall f,g \in K\left[ X \right] \setminus \{\mathbf{0}\}  \right)\left( \cont \left( fg \right) = \cont \left( f \right) \cont \left( g \right)  \right)  
	.\] 
\end{lemma} 
\begin{proof}
	TODO
\end{proof} 
\begin{lemma}
	
	Dla $f = a_0 + a_1X + \ldots + a_{n} X^{n} \in R\left[ X \right] $ mamy
	\[
	\cont \left( f \right) = 1 \iff \left(\forall p\in P \right)\left( \neg \left( p\mid a_0 \land\ldots\land p\mid a_{n}  \right)  \right)  
	.\] 
\end{lemma} 
\begin{proof}
	TODO
\end{proof} 
\begin{corollary}[Wniosek]
      Niech $f\in R\left[ X \right] \setminus R$. Wtedy następujące warunki są równoważne:
      \begin{enumerate}[label=\arabic*)]
      	\item $f$ jest nierozkładalny w $R\left[ X \right] $.
	\item $f$ jest nierozkładalny w $K\left[ X \right] $ i $\cont \left( f \right) = 1$
      \end{enumerate}
\end{corollary}
\begin{proof}
	\noindent $\left( 2 \implies 1 \right) $

	Weźmy $f_1,f_2 \in R\left[ X \right] $, takie że $f = f_1f_2$. Cel: $f_1 \in R\left[ X \right]^{\ast}  $ lub $f_2 \in R\left[ X \right]^{\ast} $.

	$f$ jest nierozkładalny w $K\left[ X \right] \implies f_1 \in K\left[ X \right]^{\ast} $ lub $f_2 \in K\left[ X \right]^{\ast} $. \\
	Załóżmy bez straty ogólności że $f_1\in K\left[ X \right]^{\ast} $. $f_1 \in R\left[ X \right] \implies f_1 = r_1 \in R\setminus \{0\} $.

	$1 = \cont \left( f \right) \overset{\text{lemat Gaussa}}{=} \cont \left( r_1f_1 \right) = \cont \left( r_1 \right) \cont \left( f_1 \right) $
	następnie

	$\cont \left( r_1 \right) \cont \left( f_1 \right) \implies \cont \left( r_1 \right) \in R^{\ast} \implies \cont \left( r_1 \right) = 1 r_1 \in R^{\ast} = R\left[ X \right] ^{\ast} $

	\noindent $\left( 1 \implies 2 \right) $

	\noindent Niech $r := \cont \left( f \right) $. Pokażemy, że $r = 1$. Wystarczy pokazać, że  $\cont \left( f \right) \in R^{\ast} $.

	Patrzymy na $f = r\cdot \tfrac{f}{r} \implies f$ nierozkładalny w $R\left[ X \right] $, $\cont \left( \frac{f}{r} \right) = 1$, $\frac{f}{r} \in  R\left[ X \right] $.

	Czyli $r\in R\left[ X \right] ^{\ast} = R^{\ast} \implies r=1$ lub $\frac{f}{r}\in R\left[ X \right] ^{\ast} = R^{\ast} \implies$ sprzeczność z założeniem, że $f \in R\left[ X \right] \setminus R$.

	Pokażemy teraz, że $f$ jest nierozkładalny w $K\left[ X \right] $. Rozważmy $g_1,g_2 \in K\left[ X \right] $, taki że $f = g_1g_2$. Cel: $g_1 \in K\left[ X \right] ^{\ast} $ lub $g_2 \in K\left[ X \right] ^{\ast} $. Niech $a_{i} := \cont \left( g_{i} \right) $, $i = 1,2$.

	Mamy

	 \[
	1 = \cont \left( f \right) = \cont \left( g_1g_2 \right) \overset{\text{lemat Gaussa}}{=} \cont \left( g_1 \right) \cont \left( g_2 \right) = a_1a_2
	.\] 
	Stąd
	\[
	f = \frac{f}{a_1a_2} = \frac{g_1}{a_1}\cdot \frac{g_2}{a_2}, \frac{g_{i}}{a_{i}} \in R\left[ X \right] 
	.\] 
	Ponieważ $f$ jest nierozkładalny w $R\left[ X \right] $ więc
	\[
	\frac{g_1}{a_1} \in R\left[ X \right] ^{\ast} \implies g_1 \in K\left[ X \right] ^{\ast} 	
	,\] 
	lub
	\[
	\frac{g_2}{a_2} \in R\left[ X \right]^{\ast} \implies g_2 \in K\left[ X \right] ^{\ast} 
	.\] 
\end{proof} 
\begin{theorem}[Gaussa]
	$R$ jest UFD $\implies R\left[ X \right] $ jest UFD
\end{theorem} 
\begin{proof}
	Niech $f \neq 0$, $f \in R\left[ X \right] \setminus R^{\ast} $. Najpierw znajdziemy rozkład $f$ na iloczyn czynników nierozkładalnych. 

	Niech $a := \cont \left( f \right) $, $f_1:= \frac{f}{a} \in R\left[ X \right] $, $\cont \left( f_1 \right) = 1 \implies f = af_1.$

	$a = a_1\cdot \ldots\cdot a_{n} $ -- rozkład $a$ na czynniki nierozkładalne (istnieje bo  $R$ jest UFD)

	$a_{i}$ jest nierozkładalny w $R \implies a_{i}$ jest nierozkładalny w $R\left[ X \right] $.

	Niech $f_1 = g_1\cdot \ldots\cdot g_{s}$ -- rozkład $f_1$ na czynniki nierozkładalne w $K\left[ X \right] $ (istnieje bo $K\left[ X \right] $ jest UFD).

	Niech $b_{i} := \cont \left( g_{i} \right) $. Mamy $1 = \cont \left( f_1 \right) \overset{\text{lemat Gaussa}}{=} b_1\cdot \ldots \cdot b_{s}$, co implikuje $f = \frac{f_1}{b_1\cdot \ldots b_{s}} = \frac{g_1}{b_1}\cdot \ldots \cdot \frac{g_{s}}{b_{s}}$, $\frac{g_{i}}{b_{i}} \in R\left[ X \right], \cont \left( \frac{g_{i}}{b_{i}} \right) = 1$. Ponadto $g_{i}$ jest nierozkładalny w $K\left[ X \right] $, $g_{i} \sim \frac{g_{1}}{b_{i}} \implies \frac{g_{i}}{b_{i}} $ jest nierozkładalny w $K\left[ X \right] $. Powyższe wynika z wniosku, wniosek implikuje $\frac{g_{i}}{b_{i}}$ jest nierozkładalne w $R\left[ X \right] $. Zatem 

	$f = af_1 = a_1\cdot \ldots\cdot a_{n} \cdot \frac{g_1}{b_1}\cdot \ldots\cdot \frac{g_{s}}{b_{s}}$ -- rozkład na iloczyn czynników nierozkładalnych w $R\left[ X \right] $. 

	Pokażemy teraz, że w $R\left[ X \right] $ każdy element nierozkładalny jest pierwszy (to wystarczy bo UFD).

	Weźmy $f\in R\left[ X \right] $ nierozkładalny oraz $g,h \in R\left[ X \right] $, takie że $f\mid gh.$ Cel: $f\mid g$ lub $f\mid h$ (w $R\left[ X \right] $ )

	\noindent Rozważmy przypadki:
	
	Przypadek $\text{deg}\left(f\right) > 0$: $f$ nierozkładalny w  $R\left[ X \right] $ więc na mocy wniosku $f$ nierozkładalny w $K\left[ X \right] $. Ponieważ $K\left[ X \right] $ jest UFD więc $f$ jest pierwszy w $K\left[ X \right] $. Ponieważ $f\mid gh$ więc $f\mid g$ lub $f\mid h$, a stąd $g = fs $ dla pewnego $s\in K\left[ X \right] $. Ale f jest nierozkładalny w $R\left[ X \right] $, zatem z wniosku wnika, że $\cont \left( f \right) = 1.$

	Ponieważ $R \ni \cont \left( g \right) = \cont \left( fs \right) \overset{\text{lemat gaussa}}{=} \cont \left( f \right) \cont \left( s \right) = \cont \left( s \right)  $ co implikuje $s\in R\left[ X \right] $. Zatem $f\mid g$ w $R\left[ X \right] $.

	Przypadek $\text{deg}\left(f\right) = 0$. $f$ jest nierozkładalny w $R\left[ X \right] \implies f$ jest nierozkładalny w $R$. Ponieważ $R$ jest UFD, więc $f$ jest pierwszy w R, a co za tym idzie $f\mid gh$, co z lematu gaussa oznacza $\cont \left( f \right) \mid \cont \left( gh \right) $ co implikuje $f\mid \cont \left( g \right) \cont \left( h \right) \implies f\mid \cont \left( g \right) $ lub $f\mid \cont \left( h \right) $, z zadania mamy wtedy $f\mid g$ lub $f\mid h$

	uwaga: $\cont \left( f \right) \sim f $ (stowarzyszenie)
\end{proof} 
\begin{corollary}
\begin{enumerate}[label=\arabic*)]
	\item $\mathbb{Z} \left[ X_1,\ldots,X_{n}  \right] $ jest UFD
	\item $K\left[ X_1,\ldots,X_{n}  \right] $ jest UFD ($K$ jest ciałem)
	\item $R\left[ X_1,\ldots,X_{n}  \right] $ jest UFD ($R$ jest UFD)
	\item Pierścień $\mathbb{Z} \left[ X_1,\ldots \right] $ jest UFD, ale nie jest noetherowski
\end{enumerate}	
\end{corollary}

\section{Kryterium Eisensteina}
Niech $R$ będzie UFD, $K = R_0$, $f = a_0 + a_1X + \ldots + a_{n}X^{n}\in R\left[ X \right] $. Ponadto niech $p$ będzie elementem nierozkładalnym w $R$. Załóżmy, że 
\[p\mid a_0,\ldots, p\mid a_{n-1}, p\nmid a_{n} , p^2\nmid a_{n} .\]
Wtedy $f$ jest nierozkładalny w  $K\left[ X \right] $.

\begin{remark}
	\[
	p\nmid a_{n} \implies a_{n} \neq 0 \implies \text{deg}(f) = n
	.\] 
	Ponadto
	\[
	p\nmid a_{n} \text{ i } p\mid a_0 \implies n>0
	.\] 
\end{remark} 
\begin{remark}
	$f$ może być nierozkładalny w $R\left[ X \right] $, bo nie mamy kontroli nad $\cont \left( f \right) $.

	\noindent (z wniosku) $f$ nierozkładalny w $R\left[ X \right] \implies \cont \left( f \right) = 1$
\end{remark} 
\begin{proof}[]Kryterium Eisensteina
	$f\sim \frac{f}{\cont \left( f \right) } \in R\left[ X \right] \text{ WLOG } \cont \left( f \right) = 1  $ (tutaj $\sim$ jest z $K\left[ X \right]$ -- nie wiem co to znaczy). Dopiero teraz mamy, że $f$ jest nierozkładalny w $K\left[ X \right] \iff f$ jest nierozkładalny w $R\left[ X \right] $ (z wniosku)

	Nie wprost, $f $ nie jest nierozkładalny (rozkładalny) w $R\left[ X \right] $ tzn. $f = gh$, dla pewnych $g,h \in R\left[ X \right] \setminus R^{\ast} $. Zauważmy, że elementy $g,h \not\in R$, bo jeśli np $g\in R$ to \[
	1 = \cont \left( f \right) = \cont \left( g \right) \cont \left( f \right) \implies g\sim \cont \left( g \right) \in R^{\ast} \implies g \in R^{\ast}  
       ,\]
       czyli sprzeczność.

Mamy $g = b_0 + b_1X + \ldots + b_{d}X^{d}$ oraz $h = c_0 + c_1X + \ldots + c_{m}X^{m}$, gdzie $b_{d} \neq 0 \neq d$ oraz $c_{m} \neq 0 \neq m$.

Następnie $f = gh \implies a_{n} = b_{d}c_{m} \neq 0$, bo $R$ jest dziedziną.

W $gh$ jest  $a_0 = b_0c_0$, ponadto z założeń $p\mid a_0$ oraz $p^2 \nmid a_0$ co nam daje $\left( p\nmid b_0 \text{ i } p\mid c_0 \right) $ lub $\left( p\mid b_0 \text{ i } p \nmid c_0 \right) $.

Następnie mamy $a_{n} = b_{d}c_{m}$ oraz załóżmy że $p\nmid a_{n} $. Wtedy $p\nmid b_{d} $ oraz $p^2 \nmid a_0$. Weźmy $r := \min \{i\le m : p\nmid c_{i}\} $. Mamy : $0<r \le m < d+m \text{(bo $d>0$)} = n$. Wtedy $a_{r} = b_0c_{r} + b_1c_{r-1} + \ldots + b_{r}c_0$ oraz $p\mid c_0 \land \ldots\land p\mid c_{r-1} \land c \nmid c_{r} \land p\nmid b_0 \implies p\nmid b_0c_0 \land p\mid b_1c_{r-1} \land \ldots \land p \mid b_{r}c_0$. Łącząc to otrzymujemy sprzecznośc z założeniem, że $p \nmid a_{r}$
\end{proof} 
\begin{eg}
	Niech $n \in \N_{>0}$ oraz $p\in \mathbb{P}$
	\begin{enumerate}[label=\arabic*)]
		\item $X^{n}-p$ jest nierozkładalny w $\mathbb{Q} \left[ X \right] $ oraz w $\mathbb{Z} \left[ X \right] $
		\item $X^{p-1} + \ldots + X + \mathbf{1}$ jest nierozkładalny w $\mathbb{Q} \left[ X \right] $ oraz w $\mathbb{Z} \left[ X \right] $
	\end{enumerate}
\end{eg} 
\subsubsection*{Kryterium związane z pierwiastkami.}
Niech $f\in K\left[ X \right] $, $\text{deg}\left( f \right) \in \{1,2,3\} $. Wtedy 
\[
f \text{ jest nierozkładalne } \iff f  \text{ nie ma pierwiastków w } K
.\] 

\section{NWD i NWW}
\begin{definition}[NWD oraz NWW]
	Niech $R$ będzie dziedziną oraz $a,b \in R$.
	\begin{enumerate}[label=\arabic*)]
		\item NWD: $c\in R$ jest \textbf{największym wspólnym dzielnikiem } (w $R$ ) $a$ i $b$, gdy
			\begin{enumerate}[label=\roman*)]
				\item $c\mid a$ oraz $c\mid b$,
				\item $\left(\forall d\in R \right)\left( d\mid a \land d\mid b \implies d\mid c \right)  $
			\end{enumerate}
		\item NWW: $c\in R$ jest \textbf{najmniejszą wspólną wielokrotnością} (w $R$ ) $a$ i $b$, gdy
			\begin{enumerate}[label=\roman*)]
				\item $a\mid c$ oraz $b\mid c$,
				\item $\left(\forall d\in R \right)\left( a\mid d \land b\mid d \implies c\mid d \right)  $
			\end{enumerate}
	\end{enumerate}
\end{definition} 

\begin{remark}
	\begin{enumerate}[label=\arabic*)]
		\item Jeśli $c$ i $c'$ są $\text{NWD}\left( a,b \right) $, $c\sim c' $. Analogicznie dla $\text{NWW}\left( a,b \right) $
		\item Jeśli $c$ jest $\text{NWD}\left( a,b \right) $ oraz $c\sim c' $ to $c'$ jest też $\text{NWD}\left( a,b \right) $. Podobnie dla $\text{NWW}\left( a,b \right) $
		\item Z 1) i 2)  $\implies$ NWD oraz NWW są wyznaczone z dosadnością do relacji $\sim  $ (jeśli w ogóle istnieją).

			NWW i NWD odnoszą się do relacji podzielności $\mid $ (????)
		\item Jeżeli $R = \mathbb{Z} $ oraz wybieramy NWD dodatni to jest on wyznaczony jednoznacznie. Podobnie NWW.
	\end{enumerate}
\end{remark} 
\begin{fact}

	Dla $c\in R$ mamy:  $c$ jest $\text{NWD}\left( a,b \right)  \iff \left( c \right) = (a)\cap (b)$
\end{fact}
\begin{proof}
	 TODO
\end{proof} 
\begin{fact}

	\[
		(c) = (a,b) \implies c \text{ jest } \text{NWD}\left( a,b \right)  
	.\] 
\end{fact}
\begin{proof}
	TODO
\end{proof} 
\begin{fact}

	Jeśli $R$ jest PID, to $c$ jest $\text{NWD}\left( a,b \right) \iff \left( c \right) = \left( a,b \right) $
\end{fact}
\begin{proof}
	TODO
\end{proof} 
\begin{fact}
Mamy:

\begin{itemize}
	\item 	$R \text{ jest UFD } \implies \left(\exists c \right)\left(c \text{ jest } \text{NWD}\left( a,b \right)\right)$
	\item $R \text{ jest UFD } \implies \left(\exists c' \right)\left( c' \text{ jest }\text{NWW}\left( a,b \right)  \right)  $
	\item $cc' \sim ab $
\end{itemize}
\end{fact}
\begin{proof}
	TODO
\end{proof} 
\begin{remark}
	\begin{enumerate}[label=\arabic*)]
		\item $R:=\mathbb{R} \left[ X,Y \right] $. Wtedy $\mathbf{1}$ jest $\text{NWD}\left( X,Y \right) $, ale $\left( \mathbf{1} \right) \neq \left( X,Y \right) $ fakt (któryś tam??) nie uogólnia się na UFD
		\item Istnieją pierścienie $R$ i $a,b \in R$ bez $\text{NWD}\left( a,b \right) $ i bez $\text{NWW}\left( a,b \right) $ 
	\end{enumerate}
\end{remark} 
\begin{proof}
	TODO
\end{proof} 
\section{Algorytm Euklidesa}

Poznamy inny sposób niż w fakcie (4?) na znajdowanie NWW i NWD w pierścieniach euklidesowych.

\noindent \textbf{Algorytm Euklidesa} 

Niech $v: R\setminus \{\mathbf{0}\} \to \N$ będzie normą euklidesową oraz $a,b\in R$ i $b\neq 0$. Szukamy $\text{NWD}\left( a,b \right) $. Przy okazji otrzymamy też $\text{NWW}\left( a,b \right) $.

\noindent \textbf{Krok 1.}: dzielimy $a$ przez $b$ z resztą.

\[
\left(\exists q_1,r_1 \right)\left( a=bq_1+r_1 \text{ oraz } \left( v\left( r_1 \right) < v\left( r_2 \right) \lor r_1 = 0 \right) \right)  
.\] 
Mamy:
\[
a = bq_1+r_1 \implies \left( a,b \right) = \left( b,r_1 \right) \overset{\text{fakt 3?}}{=} \left(\forall c\in R \right)\left( c \text{ jest }\text{NWD}\left( a,b \right) \iff c \text{ jest } \text{NWD}\left( b,r_1 \right)  \right)  
.\] 

$\text{NWD}\left( a,b \right) = \text{NWD}\left( b,r_1 \right) $-start
\begin{itemize}
	\item Jeśli $r_1 = 0 $ to $b\mid a$ i mamy $b$ jest $\text{NWD}\left( a,b \right) $.
	\item Jeśli $r_1\neq 0$ to \textbf{Krok 2.} 
\end{itemize}

\noindent \textbf{Krok 2.:} dzielimy $b$ przez $r_1$ z resztą.

\[
\left(\exists q_2,r_2 \right) \left( b = r_1q_2 + r_2  \text{ oraz } \left( v\left( r_2 \right) < v\left( r_1\right)  \lor r_2 = 0 \right) \right) 
.\] 
Jak w \textbf{Kroku 1.:} $\text{NWD}\left( b,r_1 \right) = \text{NWD}\left( r_1,r_2 \right) $ 
\begin{itemize}
	\item Jeśli $r_2 = 0$ to $r_1$ jest $\text{NWD}\left( r_1,r_2 \right) $.
	\item Jeśli $r_2\neq 0$ to \textbf{Krok 3.} 
\end{itemize}

\noindent \textbf{Krok 3.:} $r_1$ dzielimy przez $r_2$ z resztą, i tak dalej\ldots
Dostajemy ciąg: $r_1: = a$, $r_0 = ??$, $r_1,r_2,\ldots,r_{n} = 0$ 
\[
v\left( r_0 \right) > v\left( r_1 \right) > v\left( r_2 \right) > \ldots \text{ (liczby naturalne) }
.\] 
Jeśli $r_{n-1}\neq 0$ i $r_{n} = 0$, to $r_{n-1}$ jest $\text{NWD}\left( a,b \right) $

\noindent \textbf{Przykłady:}
\begin{enumerate}[label=\arabic*)]
	\item $R = \mathbb{Z} $, $v = \left| \cdot  \right| $, $a = 154, b = 42$ 
		\begin{itemize}
			\item $ \underbrace{154}_{a} = 3\cdot \underbrace{42}_{b} + \underbrace{28}_{r_1}$
			\item $\underbrace{42}_{b} = 1\cdot \underbrace{28}_{r_1}+ \underbrace{14}_{r_2}$ 
			\item $\underbrace{28}_{r_1} = 2\cdot \underbrace{14}_{r_2} + \underbrace{0}_{r_3}$
		\end{itemize}
		Zatem $14$ jest $\text{NWD}\left( a,b \right) $ 
	\item $R = \mathbb{Q} \left[ X \right] $, $v = \text{deg}$, $a = X^3-3X^2 + 3X -1$, $b= X^2-1$ 
		\begin{itemize}
			\item $a = \left( X-3 \right) \left( X^2-1 \right) + 4X - 4$
			\item $X^2 - 1 = \frac{1}{4}\left( X + 1 \right) \left( 4X - 4 \right) + 0$
		\end{itemize}
		Zatem $4X-4$ jest $\text{NWD}\left( a,b \right) $.
	\item $R = \mathbb{Z} \left[ i \right] $, $v = \left| \cdot  \right| ^2$, $a = 3-5i$, $b = 1-3i$. 
		\begin{itemize}
			\item $3-5i = 2\left( 1-3i \right) + \left( 1+i \right) $
			\item $1-3i = -\left( 1+2i \right)\left( 1+i \right) + 0  $ 
			\item (nie wiem po co ten krok): $2 = \left| 1+i \right| ^2 < \left| 1+3i \right|^2 = 10 $
		\end{itemize}
		Zatem $1+i$ jest $\text{NWD}\left( a,b \right) $
\end{enumerate}
\section{Chińskie twierdzenie o resztach.}
Niech $R$ będzie pierścieniem przemiennym z $\mathbf{1}$, $a,b \in R$, $\mathcal{I},\mathcal{J},\mathcal{J}' \unlhd R$.

Mamy 'odpowiedniości':
\begin{itemize}
	\item $\mathcal{I}\cap \mathcal{J}$ oraz NWW
	\item $\mathcal{I}+ \mathcal{J}$ oraz NWD
	\item $\left( \{cd: c\in \mathcal{I} \text{ oraz } d\in \mathcal{J}\}  \right) =: \mathcal{I}\mathcal{J}$ oraz $\cdot $ (??)
\end{itemize}
\begin{definition}
	\begin{enumerate}[label=\arabic*)]
		\item $\mathcal{I},\mathcal{J}$ są względnie pierwsze, gdy $\mathcal{I}+\mathcal{J} = R$
		\item $a,b$ są względnie pierwsze, gdy $\left( a,b \right) = R$
	\end{enumerate}
\end{definition} 
\begin{remark}
	\begin{enumerate}[label=\arabic*)]
		\item Jeśli $a,b$ są względnie pierwsze, to $\mathbf{1}$ jest $\text{NWD}\left( a,b \right) $.
		\item $R$ jest PID $\implies \left( a,b \text{ są względnie pierwsze} \iff \mathbf{1} \text{ jest } \text{NWD}\left( a,b \right)  \right)$ 
		\item $\mathbf{1}$ jest $\text{NWD}\left( X,Y \right) $ w $\mathbb{R} \left[ X,Y \right] $, ale $X$ i $Y$ nie są względnie pierwsze w $\mathbb{R} \left[ X,Y \right] $
	\end{enumerate}
\end{remark} 
\begin{lemma}
	\begin{enumerate}[label=\arabic*)]
		\item $\mathcal{I}\mathcal{J} \subseteq \mathcal{I} \cap \mathcal{J} \subseteq \mathcal{I}$,  $\mathcal{J}\subseteq \mathcal{I}+\mathcal{J}$ 
		\item $\mathcal{I}\cdot \left( \mathcal{J}+ \mathcal{J}' \right) = \mathcal{I}\mathcal{J}+ \mathcal{I}\mathcal{J}'$
		\item $\left( \mathcal{I}\mathcal{J} \right) \mathcal{J}' = \mathcal{I}\left( \mathcal{J}\mathcal{J}' \right) $
	\end{enumerate}
\end{lemma} 
\begin{proof}
	TODO
\end{proof} 
\begin{definition}[Chińskie twierdzenie o resztach]
	Niech $\mathcal{I}_{1}, \mathcal{I}_{2},\ldots, \mathcal{I}_{S}\unlhd R$, takie że $\left(\forall i\neq j \right)\left( \mathcal{I}_{i}+ \mathcal{I}_{j} =R \right)  $ (ideały są parami względnie pierwsze)

	Wtedy:
	\begin{enumerate}[label=\arabic*)]
		\item $\mathcal{I}_{1} \cap \mathcal{I}_{2} \cap \ldots \cap \mathcal{I}_{S} = \mathcal{I}_{1}\cdot \mathcal{I}_{2} \cdot \ldots\cdot * \mathcal{I}_{S}$
		\item Homomorfizm $\varphi : \mfaktor{R}{\mathcal{I}_{1}} \times \mfaktor{R}{\mathcal{I}_{2}} \times \ldots \times \mfaktor{R}{\mathcal{I}_{S}}  $ zadany wzorem
			\[
			\varphi \left( r \right) = \left( r + \mathcal{I}_{1},\ldots, r+\mathcal{I}_{S} \right) 
			,\] 
			jest ,,na'' oraz $\ker(\varphi  ) = \mathcal{I}_{1} \cap \ldots \cap \mathcal{I}_{S}$. Stąd

			\[
			\mfaktor{R}{\mathcal{I}_{1}\cap \ldots \cap \mathcal{I}_{S}} \cong \mfaktor{R}{\mathcal{I}_{1}} \times \ldots\times \mfaktor{R}{\mathcal{I}_{S}} 
			.\] 
	\end{enumerate}
\end{definition} 
\begin{proof}
	TODO
\end{proof} 
\begin{corollary}[Klasyczne CTR]
	Weźmy $n_1,n_2,\ldots,n_{\scriptscriptstyle S} \in \N_{>0}$ parami względnie pierwsze.

	Wtedy $\left(\forall a_1,\ldots,a_{\scriptscriptstyle S}\in \N \right)\left(\exists ! x_0\in \{0,1,\ldots,n_1\cdot \ldots\cdot n_{S-1}\}  \right)  $, takie że
	\[
		\left( \ast  \right) = \begin{cases}
			x_0 \equiv &a_1 \pmod{n_1} \\
			x_0 \equiv &a_2 \pmod{n_2} \\
			     &\vdots \\
			x_0 \equiv &a_{\scriptscriptstyle S} \pmod{n_{\scriptscriptstyle S}} 
		\end{cases}
	.\] 
	Ponadto zbiór wszystkich rozwiązań $\left( \ast  \right) $ to
	\[
	\{x\in \mathbb{Z} : x \equiv x_0 \pmod{n_1\cdot n_2\cdot \ldots\cdot n_{\scriptscriptstyle S}} \} = \{k\cdot n_1\cdot \ldots n_{\scriptscriptstyle S} + x_0: k\in \mathbb{Z} \} 
	.\] 
\end{corollary}
\begin{proof}
	TODO
\end{proof} 
\noindent \textbf{Przykłady:} 
\begin{enumerate}[label=\arabic*)]
	\item $n\in \N$ oraz $n=p_1^{\alpha_1}\cdot \ldots \cdot p_{s}^{\alpha_{s}}$ CTR $\implies$
		\[
		\mathbb{Z} _{n} \cong \mfaktor{\mathbb{Z} }{(n)} \overset{\cong }{\text{CTR}} \mfaktor{\mathbb{Z} }{\underbrace{p_1^{\alpha_1}}_{\mathcal{I}_{1}}} \times \ldots \times \mfaktor{\mathbb{Z} }{\underbrace{p_{s}^{\alpha_{s}}}_{\mathcal{I}_{s}}} \cong \mathbb{Z} p_1^{\alpha_1} \times \ldots \times \mathbb{Z}_{p_{s}^{\alpha _{s}}}  
		.\] 
	\item Ogólniej: Niech $R$ będzie PID   oraz $a = p_1^{\alpha_1}\cdot \ldots \cdot p_{k}^{\alpha_{k}}$ -rozkład na czynniki nierozkładalne
		\[
		\mfaktor{R}{(a)} \overset{\cong }{\text{CTR}} \mfaktor{R}{\left( p_1^{\alpha_1} \right) } \times \ldots \times \mfaktor{R}{\left( p_{k}^{\alpha_{k}} \right) } 
		.\] 
	\item Dla $R$ będącego UFD, 2) nie może zachodzić, np
		\[
		\mfaktor{\underbrace{\overbrace{\mathbb{R} }^{K}\left[ X,Y \right] }_{\text{UFD}}}{(XY)} \not\cong \mfaktor{\overbrace{\mathbb{R} }^{K}\left[ X,Y \right] }{(X)} \times \mfaktor{\overbrace{\mathbb{R} }^{K}\left[ X,Y \right] }{(Y)} \cong \overbrace{\mathbb{R} }^{K}\left[ Y \right] \times \overbrace{\mathbb{R} }^{K}\left[ X \right] 
		.\] 
		\begin{proof}
			Jako zadanie.
			Informacja: $X^3+X = X\left( X^2+1 \right) $
		\end{proof} 
	\item $\mfaktor{\mathbb{R} \left[ X \right] }{\left( X^3+X \right) } \cong \mfaktor{\mathbb{R} \left[ X \right] }{\left( X \right) } \times \mfaktor{\mathbb{R} \left[ X \right] }{\left( X^2+1 \right) } \cong \mathbb{R} \times \mathbb{C} $
	\item 
		\[\left( X \right) + \left( X+1 \right) = \mathbb{Z} \left[ X \right]\] 
		a to implikuje
		\[
		\implies \mfaktor{\mathbb{Z} \left[ X \right] }{\left( X^2+X \right) } \cong \mfaktor{\mathbb{Z} \left[ X \right] }{\left( X \right) } \times \mfaktor{\mathbb{Z} \left[ X \right] }{\left( X \right) } \times \mfaktor{\mathbb{Z} \left[ X \right] }{\left( X+1 \right) } \cong \mathbb{Z} \times \mathbb{Z} .\]

		bo $X^2+X= X\left( X+1 \right) $
\end{enumerate}
\section{Charakterystyka i pierścienie skończone.}
Niech $R$ będzie pierścieniem przemiennym z $\mathbf{1}$.

\begin{definition}[Charakterystyka.]
	$\text{char}\left( R \right) = \begin{cases}
		0, \text{ gdy }& \text{rząd}\left( \mathbf{1} \right) \text{ w } \left( R, + \right) \text{ to } \infty \\
                n, \text{ gdy }& \text{rząd}\left( \mathbf{1} \right) \text{ w } \left( R, + \right) \text{ to } n\in \N
	\end{cases}$
\end{definition} 
\noindent \textbf{Przykłady:}
\begin{enumerate}[label=\arabic*)]
	\item $\text{char}\left(\mathbb{Z} \right) = 0$.
	\item $\text{char}\left(\mathbb{Z} _{n} \right) = n$.
	\item $\text{char}\left(\mathbb{Z} _{n} \left[ X \right] \right) = n.$
	bo: $n = k\cdot l$, $k,l > 1 \implies k,l <n$. Zatem $0 = n\cdot 1 = \underbrace{ \left( k\cdot 1 \right) }_{\neq 0} \underbrace{\left( l\cdot 1 \right) }_{\neq 0}$
\end{enumerate}
\begin{fact}

	Niech $n\in \N$. Wtedy
	\[
	\text{char}\left(R\right) = n \iff \text{ istnieje monomorfizm }\varphi :\mathbb{Z} _{n} \to R \text{ zachowujący $\mathbf{1}$}  
	.\] 
	$(\mathbb{Z}_{0} = \mathbb{Z} )$
\end{fact}
\begin{proof}
	TODO
\end{proof} 
\begin{corollary}

	$R$ jest dziedziną $\implies \text{char}\left(R\right) = 0$ lub $\text{char}\left(R\right) = p \in \mathbb{P} $
\end{corollary}
\begin{proof}
	TODO
\end{proof} 
\begin{fact}

	Niech $\text{char}\left(R\right) = p\in \mathbb{P} $ oraz $a,b \in R$. Wtedy
	\[
		\left( a+b \right) ^{p} = a^{p} + b^{p}
	.\] 
\end{fact}
\begin{proof}
	TODO
\end{proof} 
\begin{corollary}

	Załóżmy, że $\text{char}\left(R\right) = p \in \mathbb{P} $. Definiujemy $F_{r}: R \to R$, $F_{r}\left( r \right) := r^{p}$.

	Wtedy $F_{r}$ jest homomorfizmem zwanym \textbf{funkcją Frobieniusa.} 
\end{corollary}
\begin{proof}
	TODO
\end{proof} 
\begin{theorem}
	$K$ jest ciałem skończonym $\implies \left| K \right| $ jest potęgą liczby pierwszej.
\end{theorem} 
\begin{proof}
	TODO
\end{proof} 
\begin{theorem}
	\begin{enumerate}[label=\arabic*)]
		\item $\left(\forall p \in \mathbb{P}  \right)\left(\forall n \in \N_{>0} \right)\left( \text{istnieje ciało mocy } p^{n} \right)   $ 
		\item $K_1, K_2$ jest ciałem skończonym, wtedy $\left| K_1 \right| = \left| K_2 \right| \iff K_1\cong K_2$
	\end{enumerate}
	$\mathbb{F} _{p^{n}}$ jest ciałem skończonym o mocy $p^{n}$ (jedyne z dokładnością do izomorfizmu)
\end{theorem} 
\begin{remark}
	$\mathbb{F} _{p} \cong \mathbb{Z} _{p}$, ale dla $n>1$ mamy $\mathbb{F}_{p^{n}} \not\cong \mathbb{Z} _{p^{n}} $
\end{remark} 
\begin{eg}
	Niech $f\in \mathbb{Z} _{p}\left[ X \right] $ - wielomiany nierozkładalne stopnia $n$. Wtedy
	\[
	\mathbb{F} _{p^{n}} \cong  \mfaktor{\mathbb{Z} _{p}\left[ X \right] }{\left( f \right) } 
	.\] 
\end{eg} 
\begin{proof}
	TODO
\end{proof} 
\begin{theorem}
	Niech $K$ będzie ciałem, $A\le K^{\ast} $, $A$ jest skończona $\implies$ $A$ jest cykliczna.
\end{theorem} 
\begin{proof}
	TODO
\end{proof} 

\section{Co to jest DVR (Discrete Valuation Ring)?}

\textbf{DVR}, czyli \textbf{pierścień dyskretnego normowania}, to szczególny rodzaj pierścienia, który stanowi „matematyczny mikroskop” pozwalający badać własności lokalne (np. zachowanie funkcji w otoczeniu punktu).

\subsection*{Definicja formalna}

\begin{definition}
DVR to \textbf{dziedzina ideałów głównych (PID)}, która ma \textbf{dokładnie jeden niezerowy ideał maksymalny}.
\end{definition}

\subsection*{Dlaczego \( R[[X]] \) jest DVR?}

W zadaniu 7 odkryłeś, że w \( R[[X]] \) jedynym elementem nierozkładalnym (z dokładnością do stowarzyszenia) jest \( X \). To ma potężne konsekwencje dla struktury tego pierścienia:

\begin{enumerate}[label=\arabic*)]
    \item \textbf{Jedyny ideał maksymalny:} Ponieważ każdy element, który nie jest wielokrotnością \( X \), jest jednostką, to ideał generowany przez \( X \), czyli \( (X) \), jest jedynym ideałem maksymalnym.
    
    \item \textbf{Struktura ideałów:} W DVR ideały są ułożone jak rosyjskie matrioszki – każdy kolejny zawiera się w poprzednim:
    \[
    \cdots \subset (X^3) \subset (X^2) \subset (X) \subset R[[X]]
    \]
    Nie ma żadnych innych ideałów! Każdy ideał jest postaci \( (X^n) \).
    
    \item \textbf{Wycena (Valuation):} Nazwa pochodzi od funkcji \( v: R[[X]] \setminus \{0\} \to \mathbb{N} \). Dla szeregu \( f \in R[[X]] \) wyceną \( v(f) \) jest najniższa potęga \( X \), która występuje w tym szeregu z niezerowym współczynnikiem (tzw. rząd szeregu):
    \[
    f = a_k X^k + a_{k+1} X^{k+1} + \cdots, \quad a_k \neq 0 \implies v(f) = k.
    \]
    \begin{itemize}
        \item \( v(f \cdot g) = v(f) + v(g) \)
        \item \( v(f + g) \geq \min\{v(f), v(g)\} \)
        \item \( v(f) = \infty \iff f = 0 \)
    \end{itemize}
\end{enumerate}

\begin{intuition}
Wyobraź sobie DVR jako pierścień, w którym możesz mierzyć „głębokość zera”. Im wyższa potęga \( X \) dzieli dany element, tym „bliżej” zera on się znajduje. 

W zwykłym pierścieniu liczb całkowitych \( \mathbb{Z} \) masz wiele różnych liczb pierwszych. W DVR masz tylko \textbf{jedną „liczbę pierwszą”} (w naszym przypadku \( X \)), która buduje całą strukturę.
\end{intuition} 

\subsection*{Inne przykłady DVR}

Poza \( R[[X]] \) najsłynniejszym przykładem są \textbf{liczby \( p \)-adyczne} \( \mathbb{Z}_p \). Jeśli weźmiesz liczbę wymierną i będziesz patrzył tylko na to, jak bardzo jest podzielna przez ustaloną liczbę pierwszą \( p \), lądujesz właśnie w świecie DVR.

\begin{eg}[Liczby \( p \)-adyczne]
Dla ustalonej liczby pierwszej \( p \), pierścień \( \mathbb{Z}_p \) jest DVR z jedynym ideałem maksymalnym \( (p) \). Każda niezerowa liczba \( p \)-adyczna ma postać \( p^n \cdot u \), gdzie \( u \) jest jednostką w \( \mathbb{Z}_p \), a \( n \geq 0 \).
\end{eg}

\subsection*{Własności DVR}

\begin{itemize}
    \item DVR jest \textbf{dziedziną lokalną} – ma dokładnie jeden ideał maksymalny.
    \item Każdy element niezerowy można zapisać jako \( \pi^n \cdot u \), gdzie \( \pi \) jest generatorem ideału maksymalnego, \( u \) jednostką, \( n \geq 0 \).
    \item DVR jest \textbf{pierścieniem Euklidesowym} względem funkcji wyceny \( v \).
    \item DVR jest \textbf{pierścieniem faktorialnym (UFD)} z dokładnie jednym elementem pierwszym (z dokładnością do stowarzyszenia).
\end{itemize}

\subsection*{Zastosowania DVR}

\begin{itemize}
    \item W geometrii algebraicznej: badanie lokalnych własności krzywych i powierzchni.
    \item W teorii liczb: konstrukcja \( p \)-adycznych uzupełnień.
    \item W analizie: badanie szeregów potęgowych i ich zbieżności.
    \item W teorii pierścieni: jako budulec bardziej skomplikowanych struktur algebraicznych.
\end{itemize}

\begin{remark}
Pierścień \( R[[X]] \) jest przykładem DVR, który powstaje naturalnie w analizie (jako pierścień szeregów formalnych) i algebraicznej geometrii (jako pierścień lokalny punktu na krzywej).
\end{remark}

